\section{Evaluierung und Ergebnisse}
In diesem Kapitel wird der Erfolg der Code-Portierung analysiert. Dabei steht insbesondere der Vergleich zwischen der ursprünglichen Python-Implementierung und der neuen TypeScript-Version im Fokus. Ein wesentliches Kriterium hierfür ist die Lesbarkeit und Nachvollziehbarkeit des Codes für Studierende.

\subsection{Lesbarkeit und syntaktischer Vergleich}
Ein zentrales Ziel der Übersetzung ist es, den Code so zu gestalten, dass dieser didaktisch wertvoll bleibt. Die Lesbarkeit von TypeScript im Vergleich zu Python hängt stark von den Vorkenntnissen der Betrachtenden ab. Für Studierende, die bereits Erfahrungen mit C-basierten Sprachen (wie C, C++ oder Java) gesammelt haben, wirkt die Syntax von TypeScript durch die explizite Typisierung und die Klammerstruktur oft vertrauter. Im Gegensatz dazu besticht Python durch seine minimalistische und pseudocode-artige Syntax, die besonders für absolute Programmieranfänger intuitiver sein kann \cite{objective-comparison}. Ein weiterer Aspekt ist, dass Python über eine Vielzahl eingebauter Datenstrukturen und Funktionen verfügt, die in TypeScript teilweise manuell nachgebildet oder über externe Bibliotheken eingebunden werden müssen.

\subsubsection{Vergleich anhand einfacher Kontrollstrukturen}
Zur Veranschaulichung wird zunächst eine einfache Funktion zur Potenzberechnung aus dem Notebook \href{https://github.com/zLeoll/TS-Translation-for-TheoInf/blob/main/Algorithms/Chapter-03/Power.ipynb}{\textit{Power.ipynb}} betrachtet.

In der ursprünglichen Python-Version (siehe Listing \ref{lst:python-power}) wird eine einfache \texttt{for}-Schleife in Kombination mit der eingebauten \texttt{range()}-Funktion verwendet, um die Multiplikation $n$-mal durchzuführen. Die Syntax ist äußerst kompakt und verzichtet in ihrer Grundform auf Typdeklarationen. Für einen direkteren Vergleich mit TypeScript wurde die Signatur hier jedoch mit Python-Typ-Hinweisen (ähnlich dem \textit{mypy}-Standard) ergänzt.
\Needspace{6\baselineskip}
\begin{minted}{python}
 def power(m: int, n: int) -> int:
     r = 1
     for i in range(n):
         r *= m
     return r
\end{minted}
\captionof{listing}{Python-Beispiel: Potenzberechnung mit Typ-Annotationen}
\label{lst:python-power}

Die äquivalente TypeScript-Implementierung (siehe Listing \ref{lst:ts-power}) zeigt die klassische, C-ähnliche \texttt{for}-Schleife. Auffällig ist hier die Wahl der Typisierung: Während die Basis \texttt{m} und das Ergebnis \texttt{r} als \texttt{bigint} deklariert sind, um das bei Potenzen rasche Wachstum ohne Präzisionsverlust abzufangen, wird der Exponent \texttt{n} als reguläre \texttt{number} getypt. Dies ist eine bewusste Designentscheidung, da abgesehen von Basiswerten wie $-1$, $0$ oder $1$, ein Exponent, der den maximalen Wert des \texttt{number}-Typs überschreitet, rechnerisch ohnehin jede verfügbare Speichergrenze sprengen würde.
\Needspace{9\baselineskip}
\begin{minted}{typescript}
	function power(m: bigint, n: number): bigint {
		let r = 1n;
		for (let i = 0; i < n; i++) {
			r *= m;
		}
		return r;
	}
\end{minted}
\captionof{listing}{TypeScript-Beispiel: Potenzberechnung}
\label{lst:ts-power}

Beim strukturellen Vergleich fällt auf, dass die Python-Variante kompakter wirkt. Dies liegt primär an der eingebauten \texttt{range}-Funktion, die in TypeScript standardmäßig fehlt. Würde man in TypeScript eine vergleichbare \texttt{range}-Utility-Funktion implementieren und mit der Array-Methode \texttt{reduce} kombinieren, ließe sich die Logik noch deutlich kürzer fassen. Da diese Hilfsfunktion jedoch erst selbst definiert werden müsste, greift die direkte Portierung hier auf die ausführlichere, aber standardkonforme \texttt{for}-Schleife zurück, welche für Entwickler mit C-Hintergrund unmittelbar verständlich ist. Beide Versionen bieten durch ihre jeweiligen Typannotationen einen klaren semantischen Mehrwert, der die Fehleranfälligkeit bei der Nutzung reduziert.


\subsubsection{Vergleich komplexerer Algorithmen}
Die strukturellen Unterschiede werden bei komplexeren Algorithmen noch deutlicher. Als Beispiel dient die Implementierung des \href{https://de.wikipedia.org/wiki/Dijkstra-Algorithmus}{Dijkstra-Algorithmus} zur Berechnung kürzester Pfade in \href{https://github.com/zLeoll/TS-Translation-for-TheoInf/blob/main/Algorithms/Chapter-11/Dijkstra.ipynb}{\textit{Dijkstra.ipynb}}.

Das Python-Listing \ref{lst:python-dijkstra} nutzt stark die eingebauten Datenstrukturen \texttt{dict} (für \texttt{Distance}) und \texttt{set} (für \texttt{Visited}). Der Code ist kompakt und profitiert von Pythons dynamischer Typisierung, was schnelle Prototypisierung ermöglicht, aber zur Laufzeit zu Fehlern führen kann, wenn unerwartete Typen übergeben werden.
\Needspace{17\baselineskip}
\begin{minted}{python}
 def shortest_path(source, Edges):
     Distance = { source: 0 }
     Visited  = { source }
     Fringe   = Set()
     Fringe.insert( (0, source) )
     while not Fringe.isEmpty():
         d, u = Fringe.pop() # get and remove smallest element
         for v, l in Edges[u]:
             dv = Distance.get(v, None)
             if dv == None or d + l < dv:
                 if dv != None:
                     Fringe.delete( (dv, v) )
                 Distance[v] = d + l
                 Fringe.insert( (d + l, v) )
         Visited.add(u)
     return Distance
\end{minted}
\captionof{listing}{Python-Beispiel: shortest\_path}
\label{lst:python-dijkstra}

Die TypeScript-Umsetzung in Listing \ref{lst:ts-dijkstra} erfordert deutlich mehr strukturellen Aufwand. Die Signaturen der Übergabeparameter (\texttt{Record<string, Array<[string, number]>>}) dokumentieren exakt die erwartete Graphenstruktur. Das Set \texttt{Visited} aus Python wurde hier durch ein \texttt{Record<string, boolean>} (ein Objekt, das als Map fungiert) ersetzt. Dies zeigt, dass TypeScript-Entwickler oft auf Objekte oder spezielle Map-Klassen zurückgreifen müssen, da die nativen Sets in TypeScript in bestimmten Anwendungsfällen weniger performant oder flexibel sind als ihre Python-Pendants.
\Needspace{31\baselineskip}
\begin{minted}{typescript}
 function shortestPath(
 source: string,
 Edges: Record<string, Array<[string, number]>>
 ): Record<string, number> {
     const Distance: Record<string, number> = { [source]: 0 };
     const visited: Record<string, boolean> = {};
     const Fringe = new AVLSet<[number, string]>();
     Fringe.insert([0, source]);
     while (!Fringe.isEmpty()) {
         const [d, u] = Fringe.pop();
         if (visited[u]) {
             continue; 
         }
         visited[u] = true;
         for (const [v, l] of Edges[u]) {
             const dv = Distance[v];
             if (dv === undefined || d + l < dv) {
                 if (dv !== undefined) {
                     Fringe.delete([dv, v]);
                 }
                 Distance[v] = d + l;
                 Fringe.insert([d + l, v]);
             }
         }
     }
     return Distance;
 }
\end{minted}
\captionof{listing}{TypeScript-Beispiel: shortestPath}
\label{lst:ts-dijkstra}

\subsubsection{Fehlende eingebaute Operatoren und Funktionen}
Ein weiterer Faktor, der die Lesbarkeit beeinflusst, ist das Fehlen bestimmter vordefinierter Operatoren in TypeScript, die in Python standardmäßig vorhanden sind. Ein relevantes Beispiel hierfür findet sich im Notebook \href{https://github.com/zLeoll/TS-Translation-for-TheoInf/blob/main/Logic/Chapter-4/06-Davis-Putnam.ipynb}{06-Davis-Putnam.ipynb}. 

In Python können \texttt{Sets} sehr elegant und lesbar mit dem Operator (\texttt{|}) vereinigt werden. Dieser Operator fasst die Menge \texttt{UsedVars} und ein neues Set, das lediglich das Element \texttt{p} enthält, in einer neuen Menge zusammen:
\Needspace{3\baselineskip}
\begin{minted}{python}
 newUsed = UsedVars | { p }
\end{minted}
\captionof{listing}{Mengenvereinigung in Python}

Da TypeScript kein Operator Overloading (Überladung von Operatoren) für komplexe Objekte unterstützt, muss dieses Verhalten methodisch nachgebildet werden. In der TypeScript-Version wurde dafür eine spezifische Methode \texttt{union} in der \texttt{RecursiveSet}-Klasse implementiert, welche die Vereinigungs-Funktionalität bereitstellt:
\Needspace{3\baselineskip}
\begin{minted}{typescript}
 const nextUsedVars = UsedVars.union(new RecursiveSet<Variable>(p));
\end{minted}
\captionof{listing}{Mengenvereinigung in TypeScript mittels \textit{.union()}}
Dieser Code erfüllt denselben Zweck, ist jedoch syntaktisch strukturierter und erfordert vom Leser das Wissen über die Existenz und Funktionsweise der spezifischen \texttt{union}-Methode auf der Klasse \texttt{RecursiveSet}.

\subsubsection{KI-generierte Code-Komplexität}
Bei der Betrachtung der Lesbarkeit muss zudem kritisch betrachtet werden, wie die verwendeten \ac{KI}-Modelle den Code übersetzt haben. Ein häufig beobachtetes Phänomen war die Generierung von unnötig komplexem oder imperativem Code, wo deklarative Ansätze lesbarer gewesen wären. 

Ein Code-Ausschnitt aus dem Notebook \href{https://github.com/zLeoll/TS-Translation-for-TheoInf/blob/main/Logic/Chapter-4/06-Davis-Putnam.ipynb}{06-Davis-Putnam.ipynb} zeigt, wie die \ac{KI} eine einfache Bedingung zunächst sehr iterativ und fehleranfällig übersetzte:
\Needspace{10\baselineskip}
\begin{minted}{typescript}
 let allUnits = true;
 for (const C of S) {
     if (C.size !== 1) {
         allUnits = false;
	     break;
     }
 }
 if (allUnits) {
     return S;
 }
\end{minted}
\captionof{listing}{Imperative (KI-generierte) Überprüfung}

Dieser Code iteriert manuell über die Menge \texttt{S}, um zu prüfen, ob alle Elemente \texttt{C} die Größe 1 haben. Diese Logik lässt sich in TypeScript durch die Nutzung von Funktionen höherer Ordnung deutlich kürzer und semantisch klarer formulieren. Erst durch manuelles Refactoring oder den expliziten Hinweis im Prompt, dass die KI wo immer möglich Konzepte der funktionalen Programmierung verwenden soll, entstand die präferierte, deklarative Variante:
\Needspace{6\baselineskip}
\begin{minted}{typescript}
 if (S.every(C => C.size === 1)) {
     return S;
 }
\end{minted}
\captionof{listing}{Deklarative (optimierte) Überprüfung in TypeScript}


Als Referenz ist hier die ursprüngliche Python-Version von \href{https://github.com/karlstroetmann/Logic/blob/master/Python/Chapter-4/06-Davis-Putnam.ipynb}{06-Davis-Putnam} zu sehen, die durch die \texttt{all()}-Funktion und List Comprehension eine ähnlich hohe Lesbarkeit wie die optimierte TypeScript-Variante aufweist:
\Needspace{5\baselineskip}
\begin{minted}{python}
 if all(len(C) == 1 for C in S):
     return S
\end{minted}
\captionof{listing}{Python-Referenz der Überprüfung}


\subsection{Laufzeit und Performance}
Neben der Lesbarkeit spielt die Ausführungsgeschwindigkeit der übersetzten Code-Basis eine entscheidende Rolle. Dies gilt insbesondere für rechenintensive Such- und Sortieralgorithmen, wie sie im Modul \textit{Theoretische Informatik II: Algorithmen und Datenstrukturen} behandelt werden. 

\subsubsection{Grundsätzliche Performance-Unterschiede}
Um die gemessenen Laufzeiten einordnen zu können, muss die zugrunde liegende Architektur der beiden Sprachen betrachtet werden. Python wird in seiner Standard-Implementierung (\href{https://en.wikipedia.org/wiki/CPython}{CPython}) zeilenweise zur Laufzeit interpretiert. Dies führt zu einem vergleichsweise hohen Overhead, da jede Instruktion bei der Ausführung in Maschinencode übersetzt werden muss. 

TypeScript hingegen wird zunächst in JavaScript kompiliert, welches dann von einer Laufzeitumgebung wie Node.js (basierend auf der V8-Engine) ausgeführt wird. Die V8-Engine nutzt einen sogenannten \textit{Just-in-Time (JIT) Compiler} \cite{v8-jit}. Dieser beobachtet den Code während der Ausführung und kompiliert häufig durchlaufene Pfade in hochoptimierten Maschinencode. Zudem hilft die strikte Typisierung von TypeScript indirekt dabei, den Code so zu schreiben, dass die V8-Engine Optimierungen wie \textit{Hidden Classes} besser anwenden kann. Folglich ist Node.js bei reinen Berechnungsaufgaben, die nicht auf externe, in C/C++ geschriebene Bibliotheken wie \href{https://numpy.org/}{NumPy} ausgelagert werden, in der Regel deutlich schneller als reines Python.

\subsubsection{Laufzeitvergleich: Merge-Sort}
Um diese theoretischen Grundlagen in der Praxis zu überprüfen, wurden Laufzeitmessungen an den portierten Algorithmen vorgenommen. Alle nachfolgenden Tests wurden auf einem Microsoft Surface 8 mit einem \textit{11th Gen Intel(R) Core(TM) i5-1135G7 @ 2.40GHz} durchgeführt.

Als exemplarisches Testobjekt dient die Implementierung des Merge-Sort-Algorithmus \href{https://github.com/zLeoll/TS-Translation-for-TheoInf/blob/main/Algorithms/Chapter-06/Merge-Sort.ipynb}{\textit{Merge-Sort.ipynb}}. Da sich der übersetzte TypeScript-Code in seiner Struktur und Logik kaum vom Python-Original unterscheidet, ist ein direkter Performance-Vergleich aussagekräftig.

Für die Messung in Python wurde die integrierte Jupyter-Notebook-\href{https://ipython.readthedocs.io/en/stable/interactive/magics.html}{\textit{Magic-Command}} \texttt{\%\%time} verwendet:
\Needspace{4\baselineskip}
\begin{minted}{python}
 %%time
 testSort(100, 2000)
\end{minted}
\caption{Time-Funktion in Python}

\textbf{Ausgabe (Python):}
\begin{verbatim}
 All tests successful!
 CPU times: total: 8.91 s
 Wall time: 18.2 s
\end{verbatim}

Für die korrekte Interpretation der Python-Ausgabe muss zwischen CPU time, der reinen Rechenzeit der Prozessorkerne, und Wall time, der real verstrichenen Zeit vom Start bis zum Ende der Zelle, unterschieden werden \cite{jupyter-walltime}. Maßgeblich für das Empfinden des Nutzers ist die Wall time von 18,2 Sekunden.

Die Messung in der TypeScript-Version erfolgte mithilfe der nativen console.time()-API. Sie misst die real verstrichene Zeit und entspricht damit der Wall time in Python:
\Needspace{5\baselineskip}
\begin{minted}{typescript}
 console.time();
 testSort(100, 2000);
 console.timeEnd();
\end{minted}

\textbf{Ausgabe (TypeScript):}
\begin{verbatim}
	All tests successful!
	default: 10.618s
\end{verbatim}

\subsubsection{Analyse der Ergebnisse (Merge-Sort)}
Der Vergleich zeigt einen signifikanten Performance-Vorteil der TypeScript-Implementierung. Mit einer Laufzeit von rund 10,6 Sekunden ist die Ausführung in Node.js/TypeScript etwa 42\,\% schneller als die Python-Variante mit 18,2 Sekunden. 

Dieser deutliche Unterschied lässt sich direkt auf die JIT-Kompilierung der V8-Engine zurückführen. Bei einem rekursiven Sortieralgorithmus wie Merge-Sort werden dieselben Funktionen und Code-Pfade tausendfach mit ähnlichen Datentypen aufgerufen. Während der Python-Interpreter diese wiederholten Aufrufe immer wieder neu verarbeiten muss, erkennt die V8-Engine diese Muster schnell und generiert optimierten Maschinencode. Dieser Performance-Gewinn ist besonders im Kontext von Lehrmaterialien vorteilhaft, da Studierende bei der Ausführung komplexer Testfälle in den Notebooks weniger Wartezeiten in Kauf nehmen müssen.

\subsubsection{\(8\)-Damen-Problem}

Das \emph{Davis--Putnam-Verfahren} ist ein klassischer Algorithmus zur Lösung von \href{https://en.wikipedia.org/wiki/Boolean_satisfiability_problem}{SAT-Problemen}, d.\,h.\ zur Entscheidung, ob eine aussagenlogische Formel in \ac{KNF} erfüllbar ist. Die Eingabe besteht dabei aus einer Menge von Klauseln. Der Algorithmus arbeitet in zwei Phasen: Zunächst wird in der \emph{Unit-Propagation} (\texttt{saturate}) jede Klausel, die nur ein einziges Literal enthält, sofort ausgewertet. Das Literal wird auf wahr gesetzt und alle betroffenen Klauseln werden entsprechend vereinfacht (\texttt{reduce}). Anschließend wählt der Algorithmus gesteuert durch eine Heuristik ein noch nicht belegtes Literal, setzt es probeweise auf wahr und ruft sich rekursiv auf. Führt ein Ast zu einem Widerspruch, wird zurückgesetzt und der andere Wahrheitswert versucht. Die \emph{Jeroslow--Wang-Heuristik} priorisiert dabei das Literal, das in möglichst vielen kurzen Klauseln vorkommt, um den Suchraum so früh wie möglich einzuschränken.

Das \href{https://de.wikipedia.org/wiki/Damenproblem}{8-Damen-Problem} wird in beiden Sprachen zunächst als SAT-Problem in \ac{KNF} kodiert und anschließend mit dem Davis-Putnam-Verfahren gelöst. Da \(n=8\) in der Praxis sehr schnell lösbar ist, wurde für eine aussagekräftigere Messung \(n=16\) gewählt, da hier die Heuristik und die zugrunde liegenden Datenstrukturen deutlich stärker beansprucht werden.\\
Der Davis--Putnam-Code erzeugt und verarbeitet in kurzer Zeit enorm viele temporäre Klauselmengen und Unit-Klauseln. In beiden Implementierungen sind genau diese ständigen Mengenoperationen und Duplikat-Prüfungen der dominierende Kostenfaktor, insbesondere in den Kernfunktionen \texttt{saturate} und \texttt{reduce}. Um diesen Architekturunterschied zu veranschaulichen, betrachten wir die \texttt{reduce}-Funktion in beiden Sprachen. Diese Funktion vereinfacht die Klauselmenge unter der Annahme, dass ein bestimmtes Literal $l$ wahr ist.

\Needspace{6\baselineskip}
\begin{minted}{python}
 def reduce(Clauses: set[Clause], l: Literal) -> set[Clause]:
     lBar = complement(l)  
     return   { C - { lBar } for C in Clauses if lBar     in C                }  \
            | { C            for C in Clauses if lBar not in C and l not in C }  \
            | { frozenset({l}) }
\end{minted}
\caption{Reduktionsschritt in Python mit Set-Comprehensions}

In Python wird die Logik extrem prägnant mit \texttt{frozenset} (Klausel) und \texttt{set} (Klauselmenge) umgesetzt. Der wesentliche Performance-Impact liegt hier in der intensiven Speicherallokation und redundanten Iterationen, die durch die kompakte Syntax verdeckt werden: Für die Reduktion muss der Interpreter drei separate Set-Comprehensions auswerten, was bedeutet, dass unter der Haube mehrfach über \texttt{Clauses} iteriert wird. Besonders teuer ist das ständige Erzeugen temporärer Objekte: Der Ausdruck \texttt{C - \{ lBar \}} zwingt Python dazu, in jedem Schleifendurchlauf ein neues Set für die Differenzbildung zu allokieren. Zudem generiert jede Comprehension eine eigene temporäre Mengen-Instanz, bevor der Vereinigungsoperator (\texttt{|}) diese Zwischenergebnisse in ein weiteres, komplett neues \texttt{set} zusammenkopiert. Dieser Overhead summiert sich bei tausenden Aufrufen in der Rekursion massiv auf.

In TypeScript wird dieselbe algorithmische Idee deutlich imperativer abgebildet. Der Code greift auf die \texttt{recursive-set}-Bibliothek zurück, da TypeScript native Set-Comprehensions für tiefgreifende Wertvergleiche nicht unterstützt:

\Needspace{22\baselineskip}
\begin{minted}{typescript}
 function reduce(Clauses: Clauses, l: Literal): Clauses {
     const lBar = complement(l);
     const result = new RecursiveSet<Clause>();
     for (const clause of Clauses) {
         let hasL = false, hasLBar = false;
         for (const lit of clause) {
             if (sameLiteral(lit, l)) hasL = true;
             if (sameLiteral(lit, lBar)) hasLBar = true;
         }
         if (hasLBar) {
             const newClause = new RecursiveSet<Literal>();
             for (const lit of clause) {
                 if (!sameLiteral(lit, lBar)) newClause.add(lit);
			 }
             result.add(newClause);
         } else if (!hasL) {
             result.add(clause);
         }
     }
     result.add(new RecursiveSet<Literal>(l));
     return result;
 }
\end{minted}
\caption{Reduktionsschritt in TypeScript mit \texttt{RecursiveSet}}

Der entscheidende Performance-Vorteil dieser textlastigeren Schreibweise liegt in der strikten Vermeidung unnötiger Allokationen. Anstatt temporäre Zwischenmengen zu erzeugen, wird initial genau eine einzige Zielmenge (\texttt{const result = new RecursiveSet<Clause>();}) allokiert, und es wird nur ein einziges Mal über \texttt{Clauses} iteriert. 
Eine neue Klausel (\texttt{new RecursiveSet<Literal>()}) wird im Speicher ausschließlich dann instanziiert, wenn das Literal \texttt{lBar} tatsächlich entfernt werden muss (\texttt{if (hasLBar)}). Bleibt eine Klausel von der Reduktion unberührt (\texttt{else if (!hasL)}), wird lediglich ihre bestehende Speicherreferenz in das neue Set (\texttt{result.add(clause)}) übernommen wodurch ein teurer Kopiervorgang völlig entfällt.

Ein zentraler Beschleuniger in der TypeScript-Version ist dabei die \texttt{recursive-set}-Datenstruktur selbst. Sie implementiert Hash-Tabellen mit \href{https://en.wikipedia.org/wiki/Open_addressing}{Open Addressing} und \href{https://en.wikipedia.org/wiki/Linear_probing}{Linear Probing} sowie eine auf Performance ausgelegte Hash-Funktion. Grundlegende Operationen wie \texttt{add}/\texttt{has} sind damit amortisiert $\mathcal{O}(1)$. 

Die im \texttt{recursive-set}-Modul formulierten Voraussetzungen (keine \texttt{NaN}/\texttt{Infinity}, keine plain objects, gute Hash-Verteilung, Unveränderlichkeit des Hashcodes nach dem Einfügen und keine Zyklen) sind in diesem SAT-Use-Case voll erfüllt, da Literale und Variablen als Strings bzw.\ \texttt{Tuple}-Objekte modelliert und die erzeugten Klauseln nach der Konstruktion nicht mehr mutiert werden. Damit greifen die effizienten \texttt{add}-Operationen im TypeScript-Code extrem schnell, was den imperativen Ansatz für das ständige Erzeugen und Prüfen von Unit-Klauseln in der Gesamtlaufzeit überlegen macht.


\subsubsection{Laufzeitvergleich: Davis--Putnam mit JW-Heuristik (\(n=16\))}
Analog zum Merge-Sort-Vergleich wurden die gleichen Methoden zum Messen der Zeit verwendet.

\Needspace{6\baselineskip}
\begin{minted}{python}
 %%time
 Solution = queens(16)
\end{minted}
\caption{Zeitmessung in Python (16-Damen)}

\textbf{Ausgabe (Python):}
\begin{verbatim}
 CPU times: total: 23.6 s
 Wall time: 24.8 s
\end{verbatim}

\Needspace{6\baselineskip}
\begin{minted}{typescript}
 console.time("queens");
 const Solution = queens(16);
 console.timeEnd("queens");
\end{minted}
\caption{Zeitmessung in TypeScript (16-Damen)}

\textbf{Ausgabe (TypeScript):}
\begin{verbatim}
	queens: 17.041s
\end{verbatim}


\subsubsection{Analyse der Ergebnisse (8/16-Damen)}
Für \(n=16\) zeigt sich ein deutlicher Performance-Vorteil der TypeScript/Node.js-Ausführung: \(17{,}041\,\mathrm{s}\) gegenüber \(24{,}8\,\mathrm{s}\) in Python, also rund \(31\,\%\) weniger \emph{Wall time}. 

Gerade beim Davis--Putnam-Verfahren treten wiederkehrende Iterationsmuster auf (z.\,B.\ in der Iteration durch Klauseln in \texttt{reduce}. Dadurch kann die V8-Engine in diesen Pfaden durch JIT-Kompilierung deutlich stärker profitieren als eine rein bytecode-interpretierte Ausführung in Python.

Im Vergleich zum Merge-Sort-Benchmark fällt der relative Performance-Abstand beim \(n\)-Damen-Problem jedoch spürbar geringer aus. Wie zuvor beschrieben, verlagern sich die Kosten bei Davis--Putnam stark in Richtung Speicherverwaltung (Aufrufe wie \texttt{add}/\texttt{has}/\texttt{union} sowie Objekterstellung). Die Laufzeit wird also weniger von engen numerischen Rechenschleifen bestimmt, in denen JIT-Compiler meist am stärksten glänzen.

Hier kommt Python ein entscheidender Architekturvorteil zugute: Bei diesem speicherintensiven Workload führt Python nicht nur reines Python aus, sondern profitiert massiv von seinen hochoptimierten, in C implementierten Hash-Tabellen für \texttt{set} und \texttt{frozenset}. Insbesondere sind zentrale Operationen wie Membership-Checks und das Hashing von \texttt{frozenset} in CPython direkt in C realisiert. Zudem wird der Hashwert eines \texttt{frozenset} nach der ersten Berechnung intern gecacht, was den Aufwand für wiederholte Hash-Berechnungen drastisch reduziert und den Performance-Nachteil des reinen Interpreters bei diesem spezifischen Problem stark abfedert.


